
\subsection{Constrained Mechanics}

Differential-algebraic equations (DAEs) are systems of equations involving both
differential and algebraic relations. In contrast to ordinary differential
equations (ODEs), a DAE does not in general determine all time derivatives
explicitly from the current state alone. Instead, some variables are linked by
constraints that must be satisfied throughout the evolution.

A key practical consequence is that not every initial condition is admissible.
The algebraic equations define a \emph{constraint manifold} on which valid
solutions must lie, and the dynamics must evolve in such a way that the
solution remains on this manifold for all time.

Constrained mechanical systems provide a canonical example. When the
configuration is described using redundant coordinates, the equations of motion
are accompanied by constraint equations enforcing geometric relations between
the coordinates. As a result, the motion is not free in the ambient coordinate
space, but is instead confined to the manifold defined by the constraints.

An important class of constraints in mechanics is that of \emph{holonomic
constraints}. A holonomic constraint depends only on the configuration
variables and may be written in the form
\begin{equation*}
    g(q) = 0,
\end{equation*}
where \(q\) denotes the generalized coordinates. Such constraints restrict the set of admissible configurations to a manifold. Velocity and acceleration are not prescribed independently, but must remain compatible with the configuration constraint so that the motion stays on the manifold over time.

This compatibility becomes clearer upon differentiating the constraint in time.
If \(g:\mathbb{R}^n \to \mathbb{R}^m\) is smooth and \(J_g(q)\) denotes its
Jacobian with respect to \(q\), then from
\begin{equation*}
    g(q(t)) = 0,
\end{equation*}
one obtains the velocity-level constraint
\begin{equation*}
    J_g(q)\,\pdv{q}{t} = 0,
\end{equation*}
and differentiating once more yields the acceleration-level constraint
\begin{equation*}
    J_g(q)\,\pdv[2]{q}{t}
    +
    \odv{}{t}\!\left(J_g(q)\right)\pdv{q}{t}
    = 0.
\end{equation*}
Thus, although a holonomic constraint is imposed at the level of configuration, it also induces hidden constraints on the allowable velocities and accelerations.

A simple example is the planar pendulum written in Cartesian coordinates. If
the bob position is given by \(q = (x,y)\), then the fixed rod length \(L\)
imposes the holonomic constraint
\begin{equation*}
    x^2 + y^2 - L^2 = 0.
\end{equation*}
The motion is therefore restricted to the circle of radius \(L\), rather than the full plane \(\mathbb{R}^2\). In this way, the pendulum equations consist not only of differential equations describing the dynamics, but also an algebraic equation enforcing the geometric constraint, and hence naturally form a DAE system.

The appearance of these hidden compatibility conditions motivates the notion of
the \emph{index} of a differential-algebraic system. Roughly speaking, the
index measures how many times the algebraic constraints must be differentiated
in time before an explicit differential system is obtained. In constrained
mechanical systems, this provides a convenient way of characterizing how
strongly the algebraic constraints are coupled to the dynamics.

For holonomically constrained systems, the constraint is imposed at the
configuration level,
\begin{equation*}
    g(q)=0.
\end{equation*}
Differentiating once gives the velocity-level compatibility condition, and differentiating a second time yields an acceleration-level relation involving \(\ddot q\). Since the equations of motion also involve accelerations, it is only at this stage that the constraint can be combined directly with the dynamics to determine the Lagrange multipliers and admissible accelerations. For this reason, the standard redundant-coordinate formulation of constrained mechanical systems is typically classified as an index-3 DAE.
